\documentclass[11pt,a4paper]{article}
\usepackage[margin=2.3cm]{geometry}
\usepackage{amsmath,amssymb,mathtools}
\emergencystretch=3em

\newcommand{\bPhi}{\boldsymbol\Phi}
\newcommand{\y}{\mathbf y}
\newcommand{\h}{\mathbf h}
\DeclareMathOperator{\vect}{vec}
\DeclareMathOperator{\diag}{diag}

\begin{document}
\begin{center}
{\LARGE The CSA Operator Is Unitary}
\end{center}

\paragraph{Setting.} Echo $\mathbf Y$ and image $\mathbf H$ are complex $M\times N$ matrices ($M$ azimuth, $N$ range samples).
$\vect(\cdot)$ stacks the columns into one vector,
\begin{equation}
  \y = \vect(\mathbf Y) = \begin{bmatrix} \mathbf Y_{:,1} \\ \vdots \\ \mathbf Y_{:,N} \end{bmatrix} \in \mathbb C^{MN},
  \qquad \h = \vect(\mathbf H) \in \mathbb C^{MN},
\end{equation}
which only rearranges entries, so $\lVert \mathbf Y \rVert_F = \lVert \y \rVert_2$.
CSA \cite{cumming2005} is linear, so it is a matrix $\bPhi$: $\h_{\mathrm{CSA}} = \bPhi\y$.

\paragraph{CSA in matrix form.} Azimuth FFT, chirp scaling, range FFT, range compression and bulk RCMC, range IFFT,
azimuth compression and residual phase compensation, azimuth IFFT:
\begin{equation}\label{eq:csa}
  \bPhi = \mathbf F_a^{-1}\mathbf P_3\,\mathbf F_r^{-1}\mathbf P_2\,\mathbf F_r\mathbf P_1\,\mathbf F_a,
\end{equation}
\begin{itemize}
  \item $\mathbf F_a = \mathbf I_N \otimes \mathbf F_M$: azimuth FFT (within each column);
        $\mathbf F_r = \mathbf F_N \otimes \mathbf I_M$: range FFT (within each row);
        $\mathbf F_K$: $K$-point DFT matrix, $\otimes$: Kronecker product.
  \item $\mathbf P_1$: chirp scaling; $\mathbf P_2$: range compression and bulk RCMC;
        $\mathbf P_3$: azimuth compression and residual phase compensation.
\end{itemize}
A product of unitary matrices is unitary, so it suffices that every factor in \eqref{eq:csa} is.

\paragraph{Phase multiplications $\mathbf P_1, \mathbf P_2, \mathbf P_3$.} Each is a Hadamard product with a unit-magnitude array,
e.g.\ $\mathbf X_{\mathrm{out}} = \mathbf E_1 \odot \mathbf X_{\mathrm{in}}$ with $[\mathbf E_1]_{m,n} = e^{j\phi_1(f_{a,m},\tau_n)}$, $\phi_1$ real.
After vectorization,
\begin{equation}
  \mathbf P_1 = \diag\bigl(\vect(\mathbf E_1)\bigr) = \diag\bigl(e^{j\phi_{1,1}}, \dots, e^{j\phi_{1,MN}}\bigr)
  \;\Rightarrow\;
  \mathbf P_1^H\mathbf P_1 = \diag\bigl(|e^{j\phi_{1,1}}|^2, \dots, |e^{j\phi_{1,MN}}|^2\bigr) = \mathbf I ;
\end{equation}
$\mathbf P_2, \mathbf P_3$ differ only in the phase function.

\paragraph{DFTs $\mathbf F_a, \mathbf F_r$.} The normalized $\mathbf F_K$ is unitary, and a Kronecker product of unitary matrices is unitary.
An unnormalized forward FFT (factor $\sqrt K$) with an inverse FFT divided by $K$ (factor $1/\sqrt K$) is also fine:
each axis has one FFT and one IFFT, so the factors cancel.

\paragraph{Conclusion.}
\begin{equation}\label{eq:unitary}
  \bPhi^H\bPhi = \bPhi\bPhi^H = \mathbf I, \qquad \bPhi^{-1} = \bPhi^H .
\end{equation}

\paragraph{Consequence: the echo-domain fidelity equals the image-domain fidelity.}
With the echo model $\mathcal A = \mathcal I^{-1}$ (inverse CSA), $\vect(\mathcal A(\mathbf H)) = \bPhi^{-1}\h$ and $\y = \bPhi^{-1}\h_{\mathrm{CSA}}$:
\begin{align}
  \bigl\lVert \mathbf Y - \mathcal A(\mathbf H) \bigr\rVert_F^2
  &= \bigl\lVert \y - \bPhi^{-1}\h \bigr\rVert_2^2
   && \text{(vectorization preserves the norm)} \notag \\
  &= \bigl\lVert \bPhi^{H}(\h_{\mathrm{CSA}} - \h) \bigr\rVert_2^2
   && \text{($\y = \bPhi^{-1}\h_{\mathrm{CSA}}$, $\bPhi^{-1} = \bPhi^H$)} \notag \\
  &= (\h_{\mathrm{CSA}} - \h)^H \bPhi\bPhi^H (\h_{\mathrm{CSA}} - \h)
   && \text{($\lVert\mathbf v\rVert_2^2 = \mathbf v^H\mathbf v$)} \notag \\
  &= \bigl\lVert \h_{\mathrm{CSA}} - \h \bigr\rVert_2^2
   = \bigl\lVert \mathbf H_{\mathrm{CSA}} - \mathbf H \bigr\rVert_F^2
   && \text{(by \eqref{eq:unitary})} .
\end{align}
So a sparse reconstruction from the echo is the same problem as one on the CSA image, which separates into per-pixel problems.
Invertibility alone is not enough: for a non-unitary $\bPhi$ the fidelity is the weighted norm
$\lVert \bPhi^{-1}(\h_{\mathrm{CSA}} - \h) \rVert_2^2$, which couples the pixels.

\begin{thebibliography}{9}
\bibitem{cumming2005}
I. G. Cumming and F. H. Wong, \emph{Digital Processing of Synthetic Aperture Radar Data: Algorithms and Implementation}. Norwood, MA, USA: Artech House, 2005.
\end{thebibliography}

\end{document}
